\subsection{半单性判据}

命题 2. --- 设 \(\varrho\) 是 G 在复希尔伯特空间 E 中的酉表示。设 \(\mathfrak{F}\) 是 \(\mathcal{M}_{c}(\mathrm{G})\) 上的滤基，并且对于紧收敛拓扑趋于测度 \(\varepsilon_{e}\)。假设存在 \(\mathrm{M} \in \mathfrak{F}\)，使得 \(\varrho(\nu)\) 对于每个 \(\nu \in \mathrm{M}\) 都是 E 的紧自同态。

那么酉表示 \(\varrho\) 是半单的，并且 G 的每个不可约酉表示在 \(\varrho\) 中的重数有限。

先证明一个引理。

引理 2. --- 设 \(\varrho_{1}\) 是 G 的一个非零酉表示，并且与 \(\varrho\) 的一个子表示同构。存在测度 \(\nu \in M\)，使得 \(\varrho_{1}(\check{\nu} * \nu)\) 是 \(\varrho_{1}\) 的表示空间上的非零紧埃尔米特自同态。特别地，存在非零实数 \(\lambda\)，使得 \(\varrho_{1}(\check{\nu} * \nu)\) 关于 \(\lambda\) 的特征空间非零。

不妨设 \(\varrho_{1}\) 是 \(\varrho\) 的子表示。令 E \(_{1}\) 为其表示空间。对于每个测度 \(\nu \in \mathrm{M}\)，自同态 \(\varrho_{1}(\nu)\) 由 \(\varrho(\nu)\) 限制到子空间得到。因此，根据假设及 III, p. 5 的命题 3，\(\varrho_{1}(\nu)\) 是紧的。

由于 \(\mathfrak{F}\) 趋于测度 \(\varepsilon_{e}\)，且收敛所在的空间是赋予紧收敛拓扑的 \(\mathcal{M}_{c}(\mathrm{G})\)，而空间 \(\mathrm{E}_{1}\) 非零，存在 \(\nu\in\mathrm{M}\)，使得 \(v=\varrho_{1}(\nu)\) 是 \(\mathrm{E}_{1}\) 的非零自同态（参见 V, p. 400, n°2）。自同态 \(u=\varrho_{1}(\check{\nu}*\nu)\) 等于 \(v^{*}\circ v\)（V, p. 401 的引理 1）；因此它非零（EVT, V, p. 39, prop. 2），是埃尔米特的且是紧的，因为 \(v\) 是紧的。

由于 \(u\) 是非零埃尔米特自同态，其谱不只含 0（I, p. 110 的例 1）。最后，由于 \(u\) 是紧的，每个 \(\lambda \in \mathrm{Sp}(u)\setminus\{0\}\) 都是 \(u\) 的特征值（III, p. 83 的命题 2）。

现在证明命题。

我们将使用 V, p. 391 的命题 7 来证明 \(\varrho\) 半单。设 \(\varrho_{1}\) 是 \(\varrho\) 的一个非零子表示，\(\mathrm{E}_{1} \subset \mathrm{E}\) 是其表示空间；我们必须证明 \(\varrho_{1}\) 包含一个不可约子表示。

取 \(\nu \in \mathrm{M}\)，使得 \(u = \varrho_{1}(\check{\nu}*\nu)\) 是 \(\varrho_{1}\) 的表示空间上的非零紧埃尔米特自同态，并取非零的 \(\lambda\)，使得 \(u\) 关于 \(\lambda\) 的特征空间 F 非零（引理 2）。空间 F 是有限维的（III, p. 90 的命题 5）。存在子表示 \(\varrho_{2}\)，它是 \(\varrho_{1}\) 在子空间 \(\mathrm{E}_{2}\)（\(\mathrm{E}_{1}\) 的子空间）上的表示，使得 \(\mathrm{E}_{2} \cap \mathrm{F}\) 非零且维数最小。取非零元素 \(x\)，它属于 \(\mathrm{E}_{2} \cap \mathrm{F}\)；令 \(\mathrm{E}_{3}\) 为 \(\varrho_{1}\) 中由 \(x\) 生成的子表示。于是 \(\mathrm{E}_{3} \subset \mathrm{E}_{2}\)，故 \(\mathrm{E}_{3} \cap \mathrm{F} = \mathrm{E}_{2} \cap \mathrm{F}\)，这是由 \(\mathrm{E}_{2} \cap \mathrm{F}\) 的维数极小性得到的。

证明表示 \(E_{3}\) 不可约。设 \(E_{4} \subset E_{3}\) 是在 G 下稳定的闭子空间。有 \(E_{2} \cap F = (E_{4} \cap F) \oplus (E_{4}^{\circ} \cap F)\)，其中 \(E_{4}^{\circ}\) 表示 \(E_{4}\) 在空间 \(E_{3}\) 中的正交补（事实上，若 \(y \in E_{2} \cap F\)，令 \(y_{4} \in E_{4}\) 为 y 在 \(E_{4}\) 上的正交投影；由于 \(E_{4}\) 和 \(E_{4}^{\circ}\) 在 u 下稳定，向量 \(u(y_{4})\) 是 \(u(y) = \lambda y\) 在 \(E_{4}\) 上的正交投影，从而 \(u(y_{4}) = \lambda y_{4}\)，这意味着 \(y_{4} \in E_{4} \cap F\)，断言得证）。于是由 \(E_{2} \cap F\) 的维数极小性，或者有 \(E_{4} \cap F = E_{2} \cap F\)，或者有 \(E_{4}^{\circ} \cap F = E_{2} \cap F\)。在第一种情形中，\(x \in E_{4}\)，从而 \(E_{4} = E_{3}\)；在第二种情形中，\(x \in E_{4}^{\circ}\)，从而 \(E_{4}^{\circ} = E_{3}\) 且 \(E_{4} = \{0\}\)。因此表示 \(E_{3}\) 不可约。

由 V, p. 391 的命题 7，得到表示 \(\varrho\) 是半单的。

设 \(\pi\) 是 G 的不可约酉表示，且其在 \(\varrho\) 中的重数非零；令 \(E_{\pi}\) 为其表示空间。存在测度 \(\nu \in M\) 和非零实数 \(\lambda\)，使得 \(u = \pi (\check{\nu} * \nu)\) 是 \(E_{\pi}\) 上的非零紧埃尔米特自同态，并且 \(u\) 关于 \(\lambda\) 的特征空间非零（引理 2）。令 F 为 \(v = \varrho (\check{\nu} * \nu)\) 关于 \(\lambda\) 的特征空间。对于每个子表示 \(E_1\)，它是 E 的与 \(\pi\) 同构的子表示；在 \(E_1\) 上由 \(v\) 限制到子空间得到的自同态与 \(u\) 等同，且 \(E_1\) 与 F 的交非零。因此，\(\pi\) 在 \(\varrho\) 中的重数小于空间 F 的维数，而该维数有限（III, p. 90 的命题 5）。

例. --- 假设 G 是单模群，例如实半单李群。设 \(\Gamma \subset G\) 是离散子群，使得商空间 \(X = \Gamma \backslash G\) 紧。群 G 通过右乘作用在 X 上。记 \(\beta\) 为 \(\Gamma\) 上的计数测度，并置 \(\widetilde{\mu} = \mu / \beta\)（INT, VII, p. 44, § 2, n° 2, 定义 1）；这是 X 上的有界 G-不变测度。

对于每个 \(f \in \mathcal{L}^2(\mathrm{X},\widetilde{\mu})\) 和每个 \(g \in \mathrm{G}\)，定义函数 \(\varrho(g)f \in \mathcal{L}^2(\mathrm{X},\widetilde{\mu})\)，使其满足 \(\varrho(g)f(x) = f(x \cdot g)\)。映射 \(\varrho(g)\) 是连续线性映射，通过取商在 \(\mathrm{L}^2(\mathrm{X},\widetilde{\mu})\) 上诱导出一个酉映射，仍记为 \(\varrho(g)\)。映射 \(\varrho\) 是 G 在 \(\mathrm{L}^2(\mathrm{X},\widetilde{\mu})\) 中的酉表示（INT, VII, p. 135, § 2, n° 5, prop. 8）。

取 \(\varphi \in \mathcal{K}(\mathrm{G})\)。对于 \(\mathrm{L}_1\) 和 \(\mathrm{L}_2\) 这两个 \(\mathrm{G}\) 的任意紧子集，\(\mathrm{T}\) 是 \(\Gamma\) 与紧集 \(\mathrm{L}_1\mathrm{Supp}(\varphi)\mathrm{L}_2^{-1}\) 的交，且该交是有限的。对于每个 \((g,h) \in \mathrm{L}_1 \times \mathrm{L}_2\)，级数 \(\sum_{\gamma \in \Gamma} \varphi(g^{-1}\gamma h)\) 与有限和 \(\sum_{\gamma \in \mathrm{T}} \varphi(g^{-1}\gamma h)\) 重合。由于 \(\mathrm{G}\) 局部紧，该级数的和 \(k_{\varphi}(g,h)\) 是 \(\mathrm{G} \times \mathrm{G}\) 上的连续函数。

取 \((g,h) \in G \times G\) 以及 \((\gamma, \eta) \in \Gamma \times \Gamma\)。有 \(k_{\varphi}(\gamma g, \eta h) = k_{\varphi}(g, h)\)，因此 \(k_{\varphi}\) 通过取商定义了 \(X \times X\) 上的连续函数，记为 \(\widetilde{k}_{\varphi}\)。又有 \(\widetilde{k}_{\varphi} \in \mathcal{L}^{2}(X \times X, \widetilde{\mu} \otimes \widetilde{\mu})\)，因为假设空间 X 是紧的。

记 \(\dot{x}\) 为 G 中元素 \(x\) 在规范投影下于 X 中的像。取 \(\varphi \in \mathcal{K}(\mathrm{G})\)。对于 \(f \in \mathcal{L}^2(\mathrm{X},\widetilde{\mu})\) 和 \(x \in \mathrm{G}\)，有

\[
\begin{array}{c} \varrho (\varphi) f (\dot {x}) = \int_ {\mathrm{G}} \varphi (g)   (\varrho (g) f) (\dot {x}) d \mu (g) = \int_ {\mathrm{G}} \varphi (g) f (\dot {x} \cdot g) d \mu (g) \\ = \int_ {\mathrm{G}} \varphi (x ^ {- 1} y) f (\dot {y}) d \mu (y) = \int_ {\Gamma \backslash \mathrm{G}} \Bigl (\sum_ {\gamma \in \Gamma} \varphi (x ^ {- 1} \gamma y) \Bigr) f (\dot {y}) d \widetilde {\mu} (\dot {y}) \end{array}
\]

（INT, VII, p. 46, § 2, no. 3, prop. 5）。由于 \(\widetilde{k}_{\varphi}\) 属于空间 \(\mathcal{L}^2 (\mathrm{X}\times \mathrm{X},\widetilde{\mu}\otimes \widetilde{\mu})\)，所以自同态 \(\varrho (\varphi)\) 在 \(\mathrm{L}^2 (\mathrm{X},\widetilde{\mu})\) 上与核为 \(\widetilde{k}_{\varphi}\) 的希尔伯特–施密特自同态重合；因此该自同态是紧的（III, p. 33 的推论 1）。

存在函数序列 \((\varphi_{n})_{n\in\mathbf{N}}\)，其元素属于 \(\mathcal{K}_{+}(\mathrm{G})\) 且积分为 1，使得 \(\varphi_{n}\cdot\mu\) 趋于测度 \(\varepsilon_{e}\)（在赋予紧收敛拓扑的 \(\mathcal{M}_{c}(\mathrm{G})\) 中）（INT, VIII, p. 139, § 2, no. 7, cor. 2）。因此命题 2 蕴含表示 \(\varrho\) 半单，且 G 的不可约酉表示在 \(\varrho\) 中的重数有限。

在 \(\varrho\) 中重数非零的 G 的不可约酉表示，称为群 G 的 \(\Gamma\)-自守表示。

